\section{Exact response identities and finite certificates}
\label{app:r42certificates}

\subsection{The inverse form from spectral cutoffs}
Let $\mathscr L$ be bounded, self-adjoint and nonnegative on a real
Hilbert space, and let $\kappa\perp\ker\mathscr L$. Write $E$ for its
spectral resolution. The response is
\[
 \chi(\kappa)=\int_{(0,\infty)}s^{-1}
                                  \,d\langle\kappa,E(s)\kappa\rangle,
\]
with $+\infty$ allowed. It satisfies
\begin{equation}\label{eq:r42spectralvariational}
 \chi(\kappa)=\sup_f
       \{2\langle\kappa,f\rangle-\langle f,\mathscr Lf\rangle\}.
\end{equation}
If the integral is finite, write
$\langle\kappa,f\rangle=
 \langle\mathscr L^{-1/2}\kappa,\mathscr L^{1/2}f\rangle$
and complete the square to obtain the upper inequality. For the
reverse inequality, use
$f_\varepsilon=\mathscr L^{-1}E([\varepsilon,\infty))\kappa$.
This is an actual Hilbert-space vector since the inverse is cut off
away from zero. Its variational value is
$\int_{[\varepsilon,\infty)}s^{-1}d\langle\kappa,E(s)\kappa\rangle$.
Monotone convergence proves equality, including the infinite case.
This also distinguishes membership in the inverse-form domain from
the stronger requirement of a Hilbert-space Poisson inverse.

\subsection{The minimum-flux characterization}
For bounded $\mathcal D:\mathcal H\to\mathcal K$ with
$\mathscr L=\mathcal D^*\mathcal D$, finite response is equivalent to
the existence of $v\in\mathcal K$ with $\mathcal D^*v=\kappa$, and
\begin{equation}\label{eq:r42minflux}
 \chi(\kappa)=\min_{\mathcal D^*v=\kappa}\|v\|^2.
\end{equation}
Indeed a flux gives
$|\langle\kappa,f\rangle|\le\|v\|\|\mathcal Df\|$ and therefore
$\chi\le\|v\|^2$. Conversely finite response defines a continuous
functional on $\operatorname{ran}\mathcal D$ by
$\ell(\mathcal Df)=\langle\kappa,f\rangle$. It is well defined
because $\kappa$ annihilates the kernel. By the sharp energy-dual
bound obtained from \eqref{eq:r42spectralvariational}, its norm is
$\sqrt\chi$. Extend to the closure and apply the Riesz theorem to
obtain $v_0\in\overline{\operatorname{ran}\mathcal D}$ with
$\mathcal D^*v_0=\kappa$ and $\|v_0\|^2=\chi$.
Every other flux differs by $\ker\mathcal D^*$, which is orthogonal
to $v_0$. Pythagoras proves minimality.

The minimizing flux need not be local. A locality theorem would be an
additional assertion. In particular a matrix correction that makes an
averaging commutator vanish algebraically may couple widely separated
coordinates.

\subsection{An exact residual-flux example}
Take the Euclidean model
\[
 L=\begin{pmatrix}2&-1\\-1&2\end{pmatrix},\quad
 \kappa=\binom10,\quad h=\binom{1/2}0.
\]
Then $L^{-1}=\frac13\left(\begin{smallmatrix}2&1\\1&2\end{smallmatrix}\right)$,
the response is $2/3$, the residual is $r=(0,1/2)^{\mathsf T}$,
and the trial value is $1/2$. Set
\[
 \mathcal Df=(f_1-f_2,f_1,f_2),\quad
 j=(-1/6,1/6,1/3).
\]
One has $\mathcal D^*\mathcal D=L$,
$\mathcal D^*j=r$ and $\|j\|^2=1/6$. Thus the upper certificate
$1/2+1/6=2/3$ is exact. The valid alternative flux $(0,0,1/2)$
gives the weaker upper bound $3/4$. These are energy-factor examples,
not identifications with a Wilson heat bath.

\subsection{A genuine finite conditional-expectation fixture}
Consider two signs $(s,t)\in\{\pm1\}^2$ in the order
$(+,+),(+,-),(-,+),(-,-)$ with weights
\[
 \nu=(3,1,1,3)/8.
\]
Let $P_1$ average $s$ at fixed $t$, and let $P_2$ average $t$ at
fixed $s$. Their matrices are
\begin{align*}
 P_1&=\begin{pmatrix}
 3/4&0&1/4&0\\0&1/4&0&3/4\\
 3/4&0&1/4&0\\0&1/4&0&3/4
 \end{pmatrix},\\
 P_2&=\begin{pmatrix}
 3/4&1/4&0&0\\3/4&1/4&0&0\\
 0&0&1/4&3/4\\0&0&1/4&3/4
 \end{pmatrix}.
\end{align*}
They are projections self-adjoint for the weighted inner product.
The functions $s+t$, $s-t$ and $st-1/2$, together with the constant,
form an orthogonal eigenbasis for $L=2I-P_1-P_2$. Their eigenvalues
are $1/2$, $3/2$, $2$ and $0$, respectively. This verifies directly
that the true gap is $1/2$. The density-ratio comparison with the
uniform product measure gives the valid lower bound $\lambda_0=1/3$.
Thus the constructive residual rate of Theorem~\ref{thm:r42neumann}
uses $\rho=5/6$, while the true centred norm of
$T=(P_1+P_2)/2$ is $3/4$.

For $F=t$ and $m=1$, $h^{[1]}=t/2$ and
$r^{[1]}=Tt=s/4+t/2$. The two sources $s$ and $t$ have different
coordinate supports, so
\[
 \sum_i\langle Q_is,Q_ih^{[1]}\rangle=0,
 \qquad \langle s,r^{[1]}\rangle=1/2=\Cov(s,t).
\]
This checks the exact residual identity with the off-diagonal pair
retained. It does not assert that its small finite law is a
nonabelian Wilson model.

\subsection{Gaussian transport and entropy checks}
For the conditional Gaussian family with covariance $\tau C(x)^{-1}$
and mean $m(x)$, assume $C$ and $m$ are continuously differentiable,
$C>0$, and take initially smooth test functions of compact support.
The symmetric solution
\[
 CK_i+K_iC=-\partial_i C
\]
gives a velocity $v_i=m_i+K_i(y-m)$ satisfying
\[
 \partial_i\E_x F=\E_x\partial_iF+\E_x(v_i\cdot\nabla_yF).
\]
In an orthonormal eigenbasis of $C$, its entries are
$(K_i)_{ab}=-(C_i)_{ab}/(\lambda_a+\lambda_b)$, which proves existence
and uniqueness without assuming that $C$ commutes with $C_i$.
Gaussian integration by parts proves the derivative identity: its score is
\[
 \tfrac12\Tr(C^{-1}C_i)
 -\tfrac1{2\tau}(y-m)^{\mathsf T}C_i(y-m)
 +\tau^{-1}m_i^{\mathsf T}C(y-m).
\]
For a trial symmetric $H_i$ define
$R_i=C_i+CH_i+H_iC$. If $C\ge cI>0$, the error
$D_i=H_i-K_i$ satisfies
\[
 \|D_i\|_F\le\frac{\|R_i\|_F}{2c},\qquad
 \E\|D_i(y-m)\|^2\le\frac{\tau\|R_i\|_F^2}{4c^3}.
\]
Pair $CD_i+D_iC=R_i$ with $D_i$ in Frobenius norm, then use
$\tau C^{-1}\le(\tau/c)I$. No simultaneous diagonalization is used.

At $\tau=1$, $C=\operatorname{diag}(1,2)$ and
$C_i=\left(\begin{smallmatrix}0&1\\1&0\end{smallmatrix}\right)$,
the optimum symmetric matrix has off-diagonal entries $-1/3$.
Its velocity cost is $1/6$. The exact nonsymmetric transport
$-C^{-1}C_i/2$ costs $3/16$. A symmetric trial with off-diagonal
entries $-1/4$ has residual off-diagonal $1/4$ and inverse-energy
residual response $1/96$. These rational identities are checked by
the ancillary program.

Here optimality has an exact affine meaning. Set $\Sigma=\tau C^{-1}$.
An affine covariance velocity $A_i$ is admissible when
$A_i\Sigma+\Sigma A_i^{\mathsf T}=\partial_i\Sigma$.
The symmetric Sylvester solution $K_i$ is admissible. For
$D_i=A_i-K_i$ one has $D_i\Sigma+\Sigma D_i^{\mathsf T}=0$.
Expanding the cost and taking the trace gives
\[
 \Tr(A_i\Sigma A_i^{\mathsf T})
 =\Tr(K_i\Sigma K_i)+\Tr(D_i\Sigma D_i^{\mathsf T}),
\]
since the cross term is
$\Tr[K_i(D_i\Sigma+\Sigma D_i^{\mathsf T})]=0$.
The positive-definite covariance makes equality possible only when
$D_i=0$. This proves minimum mean-square velocity among admissible
affine transports; it asserts neither locality nor a Wilson transport.

Retain all coarse directions in the Gram matrix
\[
 \Gamma_{ij}=\E(v_i\cdot v_j)
       =m_i\cdot m_j+\tau\Tr(K_iK_jC^{-1}).
\]
For any parameter vector $a$, the derivative identity and
Cauchy--Schwarz give
\begin{align*}
 &|a\cdot(\nabla_x\E_xF-\E_x\nabla_xF)|^2\\
 &\hspace{8mm}\le(a^{\mathsf T}\Gamma a)\E_x|\nabla_yF|^2.
\end{align*}
Taking the supremum over unit $a$ gives the same bound with
$\|\Gamma\|_{\rm op}$ for the squared vector norm. Replacing the
operator norm by its trace can add an avoidable factor equal to the
number of directions. If $\mathcal R(a)=\sum_i a_iR_i$ has operator
norm $r$ into Frobenius matrices and $\Gamma_H$ is the Gram matrix of
the trial velocities, the preceding residual bound applied to each
combined direction gives
\[
 \sqrt{\|\Gamma\|_{\rm op}}
 \le\sqrt{\|\Gamma_H\|_{\rm op}}+\frac{\sqrt\tau\,r}{2c^{3/2}}.
\]
This follows from the triangle inequality for the maps from the
parameter space to $L^2(\nu_x;\R^m)$. Its application over a family
still requires the indicated precision and Gram bounds uniformly.

For clarity about what such Gaussian tools do not infer, consider
\[
 V(x,y)=x^2/32+\tfrac{c_0}2(1+x^2)y^2,
 \qquad c_0=487/5000.
\]
The fine minimizer is zero and its profiled curvature is $1/16$,
while the fine Hessian is at least $c_0$. Exact integration at
$\tau=1$ gives, up to a constant,
\[
 W(x)=x^2/32+\tfrac12\log(1+x^2),\qquad W''(\sqrt3)=-1/16.
\]
Equivalently at that point the expected direct curvature is $5/16$
and the force variance is $3/8$. This proves that a positive fine
Hessian and a positive minimized coarse curvature do not imply
positive curvature of the integrated free energy.

\subsection{A positive variance gap despite the entropy obstruction}
The negative curvature in the preceding example is not an obstruction
to every variance inequality. It is possible to close that question
for the specified model, with a completely explicit constant.

\begin{proposition}\label{prop:r42nonconvex}
Let $\mu(dx\,dy)$ be the normalized probability measure proportional
to $e^{-V(x,y)}dx\,dy$ for the two-variable potential just displayed.
For its weighted Sobolev form domain,
\begin{equation}\label{eq:r42nonconvexgap}
 \Var_\mu(F)\le\frac{2696920}{4383}
                \E_\mu(|F_x|^2+|F_y|^2).
\end{equation}
This is a variational gap for this probability model, with the
gradient Dirichlet form appearing in \eqref{eq:r42nonconvexgap}.
\end{proposition}
\begin{proof}
The coarse law has density proportional to
$e^{-x^2/32}/\sqrt{1+x^2}$. Write
\[
 r(x)=\frac{x^2}{2(1+x^2)},\qquad U(x)=W(x)-r(x).
\]
Then $0\le r<1/2$ and
\[
 U''(x)=1/16+\frac{x^2(3-x^2)}{(1+x^2)^3}\ge1/144.
\]
For the last inequality put $u=x^2$. After multiplication by
$18(1+u)^3$ the claim is
$p(u)=u^3-15u^2+57u+1\ge0$ for $u\ge0$.
For $0\le u\le1$, $p(u)\ge u^3+42u+1$.
For $1\le u\le7$,
$p(u)=(u-7)^2(u-1)+50-6u\ge8$.
For $u\ge7$,
$p(u)=(u-8)^2(u+1)+9(u-7)\ge0$.

A one-dimensional density proportional to $e^{-U}$ with
$U''\ge\alpha>0$ has variance constant at most $1/\alpha$.
Here is a direct justification. For centred smooth compactly
supported $f$, let $\rho$ be the normalized density and choose
\[
 g'(x)=-\rho(x)^{-1}\int_{-\infty}^x f(s)\rho(s)\,ds.
\]
The centering makes $g'$ compactly supported, and
$Lg=-g''+U'g'=f$. Integration by parts gives
$\E f^2=\E f'g'$ and
\[
 \E(Lg)^2=\E(g'')^2+\E[U''(g')^2]\ge\alpha\E(g')^2.
\]
Cauchy--Schwarz proves $\E f^2\le\alpha^{-1}\E(f')^2$.
Truncation, a compactly supported mean correction and local smoothing
extend this to the centred form domain. Strong convexity ensures the
needed integrability. Comparing $e^{-W}$ to $e^{-U}$ costs at most
$e^{\operatorname{osc}r}=\sqrt e$, so the coarse variance constant is
at most $144\sqrt e<240$. One rational check of the last inequality
is $e\le65/24+1/100=1631/600<25/9$.

The conditional law of $y$ at fixed $x$ is Gaussian of variance
$v(x)=1/[c_0(1+x^2)]$. It has variance constant at most $1/c_0$.
For $g(x)=\E[F\mid x]$, Gaussian integration by parts gives
\[
 g'=\E[F_x\mid x]-\frac{x}{1+x^2}\E[yF_y\mid x].
\]
Indeed the score of this conditional density is
$-c_0x(y^2-v)$ and
$\Cov(y^2,F)=v\E(yF_y)$.
Conditional Cauchy--Schwarz and
$\sup_{u\ge0}u/(1+u)^3=4/27$ give
\[
 \|g'\|_2\le\sqrt{\E_\mu|F_x|^2}
        +\sqrt{\frac4{27c_0}}\sqrt{\E_\mu|F_y|^2}.
\]
Total variance is consequently bounded by
\[
 c_0^{-1}E_y+240\bigl(\sqrt{E_x}
                     +\sqrt{4/(27c_0)}\sqrt{E_y}\bigr)^2,
\]
where $E_x=\E_\mu|F_x|^2$ and $E_y=\E_\mu|F_y|^2$.
The positive two-by-two comparison matrix has trace
\[
 c_0^{-1}+240(1+4/(27c_0))=2696920/4383.
\]
Its largest eigenvalue is at most its trace. This proves
\eqref{eq:r42nonconvexgap}, first for smooth test functions and then
by closure. Normalizability follows directly from the displayed
Gaussian-decaying coarse density and its conditional Gaussian law.
\end{proof}

For a product of finitely many copies of this specified model,
tensorization preserves the same constant. It does not do so for an
interacting Wilson family without another dependence estimate.
Proposition~\ref{prop:r42nonconvex} shows that a failed convexity
criterion can be bypassed by a proved two-level comparison in an
explicit model. It does not supply that comparison for the Wilson
measure or identify this gradient gap with a physical mass.

\subsection{The limits of matrix-order certificates}
A matrix upper bound cannot be interpreted entrywise. For example
\[
 X=\begin{pmatrix}1&1/2\\1/2&1\end{pmatrix}\le2I,
\]
but $X_{12}\ne0=(2I)_{12}$. A bounded positive operator norm does not
give a volume-uniform absolute row bound either. For an integer
$k\ge1$ let $v=(1,1/k,\ldots,1/k)$ have $k^2+1$ entries and set
$P=vv^{\mathsf T}/2$. Then $\|v\|^2=2$, so $P$ is a rank-one
orthogonal projection with norm one. Its first absolute row sum is
$(k+1)/2$, which diverges. All its entries are rational.

\subsection{Checks and the remaining quantitative problem}
The file \texttt{anc/check-response.py} independently checks the
finite rational examples, the conditional projection identities,
the finite Poisson residuals and the coefficient arithmetic. Its
output is \texttt{anc/response-checks.json}. These checks validate
specified finite identities; the general Hilbert-space arguments are
the written proofs above and in Secs.~\ref{ssec:r42off}--\ref{ssec:r42flux}.
No new Lean kernel replay or Wilson Monte Carlo experiment is claimed.

To obtain (FC), the unresolved task is to construct trials or fluxes
for the actual constrained Wilson forces with a sufficiently small
off-diagonal bound, uniformly over the required coarse data, volumes
and couplings. The explicit finite-volume rate
\eqref{eq:r42residualrate} proves existence and makes its cost visible,
but does not meet that uniform target. The scale-by-scale effective
action and the nontrivial continuum construction remain separate
requirements after a successful local certificate.
